\magnification=\magstep1
\input notesmac.tex
\advance\vsize by 1cm
\nopagenumbers
\centerline{\titlefont Tech and Tools, Lab 4}
\bigskip
\bigskip
\line{Dr.~WANG Jian-Sheng \hfill for 28/30 Sep 1999}
%----------------------------------------------------------------------- 
\bigskip
\problem The three-dimensional laplacian operator in cartesian coordinates
$(x,y,z)$ is given as 
$$ \nabla^2  = {\partial^2 \over \partial x^2 } + 
   {\partial^2 \over \partial y^2 } + 
   {\partial^2 \over \partial z^2 }.  $$
Using Mathematica, find the same operator in spherical polar coordinates
$(r, \theta, \phi)$, related to cartesian coordinates by
$$\eqalign{ x & = r \sin\theta \cos\phi, \cr
            y & = r \sin\theta \sin\phi, \cr
            z & = r \cos\theta. \cr}
$$
This involves viewing $f$ (in $\nabla^2 f$) as a function of
$r$, $\theta$, and $\phi$, using the chain rule to compute partial
derivatives.  [N.B. we want to find $\hat \nabla^2$ such
that $\nabla^2 f(x,y,z) = \hat \nabla^2 \hat f(r,\theta,
\phi)$ where $\hat f(r,\theta, \phi) = f\Bigl(x(r,\theta,\phi), 
y(r,\theta,\phi),z(r,\theta,\phi)\Bigr)$.]

\bigskip
\problem Write a one-liner function sort (by the use of rule-based
programming) which sorts a list of numbers in ascending order.
For example {\tt sort[\char`{8,4,2,3,2,4,2\char`}]} produces 
{\tt \char`{2,2,2,3,4,4,8\char`}}.

\bigskip
\problem Develop a set of rules for Laplace transformation.  
The Laplace transform on a function is defined by
$$ F(s) = \int_0^\infty e^{-st} f(t)\,dt.$$
The Laplace transform has the following properties:
\item{} (a) $\int_0^\infty e^{-st}[ f(t) + g(t)]\,dt = F(s) + G(s)$ (linearity).
\item{} (b) $\int_0^\infty e^{-st} c f(t)\,dt = c F(s)$ (linearity).
\item{} (c) $\int_0^\infty e^{-st} c\,dt = c/s$, where $c$ is a constant. 
\item{} (d) $\int_0^\infty e^{-st} t^n dt = n!/s^{n+1}$, where $n>0$ is an integer .
\item{} (e) $\int_0^\infty e^{-st} g(t) t^n dt = (-1)^n {d^n G(s)/ds^n}$, where $n>0$ is an integer .
\item{} (f) $\int_0^\infty e^{-st} g(t) e^{b+ct} dt = e^b G(s-c)$, where $b$ 
and $c$ are constants

\item{} Test your implementation by evaluating the Laplace transform of
$t^2 e^{c t + b}$, where $b$ and $c$ are constants.  Also compare your
implementation with built-in function {\tt LaplaceTransform[ ...]}.

\bye
